
\section{Validation Strategy and Reference Cases}
Before applying any CFD solver to a new or unfamiliar geometry, it is essential to first verify that the solver produces reliable results for a known case. Without this step, there is no way to judge whether the predictions for the actual target geometry are physically meaningful or affected by errors in the simulation setup. This is especially important when using OpenFOAM, because unlike commercial software such as ANSYS Fluent, it does not provide a graphical interface that guides the user through the setup process. Every setting must be configured manually through text files, which increases the risk of mistakes in boundary conditions, solver parameters, or mesh configuration. A structured validation study helps to identify and eliminate such errors before moving to the actual CAFA drone simulation.

The validation approach used in this thesis is based on a three-way comparison: OpenFOAM simulation results are compared against both wind tunnel experimental data and ANSYS Fluent predictions for the same geometry and flow conditions. This approach has two main advantages. First, comparison with experimental data shows how close the simulation is to real physical behaviour. Second, comparison with ANSYS Fluent results allows a direct evaluation of OpenFOAM against an established and widely trusted commercial solver. If OpenFOAM produces results that are close to both the experiment and ANSYS Fluent, it can be considered sufficiently accurate for the subsequent drone analysis.

The reference case selected for this validation is the SSAM-Gen5 model (Sydney Standard Aerodynamic Model -- Generation~5), a generic delta wing aircraft geometry. This model was chosen for several reasons. First, it is a delta wing configuration with sharp leading edges that produces complex vortex-dominated flow at moderate and high angles of attack \cite{Giannelis2023SSAMGen5}. This makes it a challenging test case for any CFD solver, because capturing vortex formation, interaction, and breakdown correctly requires good mesh quality and appropriate solver settings. If OpenFOAM can handle this level of flow complexity, it can be expected to perform well for the simpler flow around the CAFA Tech drone. Second, the SSAM-Gen5 model length is $0.75~\mathrm{m}$, which is comparable in scale to the CAFA Tech drone. This means that the Reynolds number range and mesh resolution requirements are similar, making the validation directly relevant to the final application. Third, all geometry files, experimental data, and computational meshes are publicly available through the Zenodo repository \cite{Giannelis2023SSAMGen5}, which ensures full transparency and reproducibility of the validation process.

The experimental reference data were obtained in subsonic wind tunnel tests at the University of Sydney at a freestream velocity of $20~\mathrm{m/s}$, corresponding to a Reynolds number of approximately $3 \times 10^{5}$ based on the mean aerodynamic chord \cite{Giannelis2023SSAMGen5}. Force and moment measurements were recorded across a range of angles of attack from $-5^{\circ}$ to $30^{\circ}$ \cite{Giannelis2023SSAMGen5}. The experimental results confirmed that the SSAM-Gen5 exhibits aerodynamic behaviour typical of slender delta wing aircraft, including monotonically increasing lift without abrupt stall and a pitching moment break near $12^{\circ}$ angle of attack \cite{Giannelis2023SSAMGen5}.

The ANSYS Fluent reference results were taken from the same publication by Giannelis et al.\ \cite{Giannelis2023SSAMGen5}. These simulations used the steady, incompressible, pressure-based solver with second-order spatial discretisation \cite{Giannelis2023SSAMGen5}. Several turbulence models were tested, including Spalart--Allmaras (SA), $k$--$\omega$ SST, and Realisable $k$--$\varepsilon$, both with and without curvature correction \cite{Giannelis2023SSAMGen5}. The results showed that the turbulence model choice has a strong effect on the predicted aerodynamic coefficients at high angles, with the SA model providing the best correlation with the experimental data overall \cite{Giannelis2023SSAMGen5}. For the validation in this thesis, the OpenFOAM results are compared against the ANSYS Fluent predictions obtained with both the $k$--$\omega$ SST and Spalart--Allmaras models, as these are the same two turbulence models used in the OpenFOAM simulations.

To ensure a fair comparison, all three data sources (wind tunnel, ANSYS Fluent, and OpenFOAM) use the same model geometry, the same freestream velocity of $20~\mathrm{m/s}$, and the same angle of attack range. The same reference area and reference length are used for computing the aerodynamic coefficients. The only differences lie in the mesh generation approach (\texttt{snappyHexMesh} in OpenFOAM versus Pointwise in ANSYS Fluent) and the internal numerical implementation of the solvers. These are precisely the differences that the validation study is designed to evaluate.

\section{Geometry: SSAM-Gen5 Delta-Wing Configuration}
In this section the geometry and experimental setup from Giannelis et al. \cite{Giannelis2023SSAMGen5} is summarized.
The geometry used for validation is the SSAM-Gen5 model, a generic aircraft designed to represent the key shape features of modern fifth-generation high-performance platforms. The overall shape is based on the general layout of the F-22 Raptor, but with significant simplifications to keep the geometry easy to manufacture and to model in CFD.

The full-scale vehicle length is $19~\mathrm{m}$, and the wind tunnel model is scaled to $0.75~\mathrm{m}$, which corresponds to approximately a 1:25 scale. The main reference values for the scaled model are summarised in Table~\ref{tab:ssam_ref_values}.

\begin{table}[H]
\centering
\caption{Reference values for the SSAM-Gen5 model at $0.75~\mathrm{m}$ scale \cite{Giannelis2023SSAMGen5}.}
\label{tab:ssam_ref_values}
\begin{tabular}{lcc}
\hline
\textbf{Parameter} & \textbf{Value} & \textbf{Unit} \\
\hline
Reference area          & 0.1091 & $\mathrm{m^2}$ \\
Mean aerodynamic chord  & 0.2265 & $\mathrm{m}$ \\
Span                    & 0.5350 & $\mathrm{m}$ \\
\hline
\end{tabular}
\end{table}

The wing profiles are four-percent-thick biconvex sections with sharp leading edges. No camber or spanwise twist is applied to the wings, which keeps the geometry as simple as possible while still producing realistic vortex-dominated flow. The sweep angle of the main delta wing and the horizontal stabiliser is approximately $42^{\circ}$, while the twin vertical stabilisers are swept back about $25^{\circ}$ and canted outwards at $27.5^{\circ}$. The aircraft includes engine intake openings, but these are simplified and sealed in the wind tunnel model. Thrust nozzles at the rear of the fuselage are retained in the geometry.

The sharp leading edges are an important feature of this geometry. Unlike wings with rounded leading edges, sharp edges force flow separation to occur at a fixed location along the entire leading edge, which leads to the formation of strong vortices over the upper surface of the wing. These vortices produce additional lift at moderate and high angles of attack, known as vortex lift. At the same time, the vortical flow structures make the aerodynamic predictions sensitive to mesh quality, turbulence model choice, and solver settings. This is what makes the SSAM-Gen5 a demanding and therefore useful validation case.

The wind tunnel model was fabricated from 3D-printed PLA parts, which were assembled, sanded, and coated with resin to provide a smooth surface finish. The model was designed with removable wings, fins, and horizontal stabilisers to allow testing of different configurations.

The STL geometry file used in the OpenFOAM simulations was downloaded directly from the publicly available Zenodo repository. This ensures that the exact same geometry is used in both the OpenFOAM and ANSYS Fluent simulations, eliminating any differences that could arise from geometry reconstruction or approximation.

\section{Computational Domain and Mesh Generation}
The computational domain for the simulations was defined as a rectangular box around the delta-wing model. The domain size was selected to reduce the influence of the outer boundaries on the flow around the wing, even at high angles of attack (above $20^\circ$). A sufficiently long downstream distance also helps to avoid reverse flow at the outlet and allows the wake behind the aircraft to develop inside the domain.

The delta-wing model length is $0.75~\mathrm{m}$. The domain limits were defined in Cartesian coordinates as:
\begin{itemize}
    \item $x$-direction: from $-3.0~\mathrm{m}$ to $7.5~\mathrm{m}$,
    \item $y$-direction: from $0.0~\mathrm{m}$ to $3.0~\mathrm{m}$,
    \item $z$-direction: from $-3.0~\mathrm{m}$ to $3.0~\mathrm{m}$.
\end{itemize}
Only half of the geometry was simulated by applying a symmetry plane at $y = 0$, which reduces the computational cost while remaining valid for a symmetric geometry and flow setup.

The background mesh was generated using the \texttt{blockMesh} utility, creating a three-dimensional hexahedral block. Boundary patches were assigned to each outer face of the domain, including \texttt{inlet} (front patch), \texttt{outlet} (rear patch), \texttt{lowerWall}, \texttt{upperWall}, and \texttt{sideWall}. A \texttt{symmetryPlane} boundary condition was defined to enable the half-domain simulation.

After the background mesh and boundaries were created, the \texttt{surfaceFeatureExtract} utility was used to extract feature edges from the delta-wing \texttt{.stl} geometry. This step helps preserve sharp geometric details (for example, leading and trailing edges) during the refinement and snapping stages. The main parameter of this utility is \texttt{includedAngle}, which identifies feature edges based on the angle between adjacent surface normals. For the delta-wing geometry, an \texttt{includedAngle} of $176^\circ$ was selected after several trials to capture sharp wing edges reliably.

Mesh generation was then performed using the \texttt{snappyHexMesh} utility, which produces a predominantly hexahedral mesh fitted to the surface geometry. The initial setup for this step was based on the official OpenFOAM \texttt{motorBike} tutorial and then adjusted through multiple attempts to obtain acceptable mesh quality around the wing.

To improve resolution in regions with strong flow gradients, two refinement zones were added using refinement boxes:
\begin{itemize}
    \item a \textbf{near-field} refinement box, slightly larger than the delta-wing size, to resolve vortices and strong gradients close to the aircraft,
    \item a \textbf{far-field} refinement box, larger than the near-field zone and extended downstream, to capture the wake behaviour after the flow passes the aircraft.
\end{itemize}
These refinement regions increase mesh density only where needed, keeping the total number of cells at a reasonable level.

\subsubsection{Boundary Layer Mesh}
Boundary layers are one of the most important parts of mesh generation in external aerodynamics, because aerodynamic forces strongly depend on the near-wall velocity gradient. In addition, the layer region often contains a large portion of the final mesh cells.

Two final meshes were generated for comparison:
\begin{itemize}
    \item \textbf{Mesh 1:} target $y^+ \approx 4$,
    \item \textbf{Mesh 2:} target $y^+ \approx 5$.
\end{itemize}
The first cell height was estimated using a $y^+$ calculator, and at least 10 prism layers were added on the wing surface. With 10 layers, the layer growth ratio was kept smaller than 1.15. Growth ratios larger than 1.2 are generally considered poor practice because the transition between layers becomes too aggressive, which can reduce accuracy in the boundary layer and near separation regions.

During the layer generation stage, a practical issue was observed: \texttt{snappyHexMesh} can produce low-quality or collapsed layers when trying to generate more than about 5 layers in a single run. To avoid this, the layer addition was performed in stages:
\begin{enumerate}
    \item \texttt{snappyHexMesh} was first run with \texttt{castellatedMesh} and \texttt{snap} enabled, while \texttt{addLayers} was disabled, producing a clean mesh without layers.
    \item The mesh from the first step was then used as the initial mesh, and \texttt{snappyHexMesh} was run again to add 5 layers.
    \item After verifying layer quality, \texttt{snappyHexMesh} was run once more to add the remaining 5 layers. During this step, the \texttt{finalLayerThickness} parameter was reduced by a factor of two, because half of the total layer thickness already existed from the first layer run.
\end{enumerate}
This staged method produced stable and good-quality layers for both $y^+$ targets.

\subsubsection{Mesh Quality Check}
At the end of the meshing process, the \texttt{checkMesh} utility was used to evaluate mesh quality parameters such as non-orthogonality, skewness, and overall mesh consistency. Both meshes contained a small number of skewed faces, with a maximum skewness value of approximately 11. For large meshes, a small number of locally skewed cells is typically acceptable; however, these regions should be monitored because poor-quality cells can affect solver stability and accuracy.

Finally, the main mesh parameters are summarised in Table~\ref{tab:mesh_parameters} (to be added), including the total number of cells, number of prism layers, first-layer thickness, estimated $y^+$ range for each mesh, and key mesh quality metrics.

% Example placeholder table (fill with your real values):
% \begin{table}[h!]
% \centering
% \caption{Final mesh parameters for both meshes.}
% \label{tab:mesh_parameters}
% \begin{tabular}{lcc}
% \hline
% Parameter & Mesh 1 ($y^+\approx 4$) & Mesh 2 ($y^+\approx 5$) \\
% \hline
% Total cells &  &  \\
% Prism layers & 10 & 10 \\
% Growth ratio & 1.11 & 1.11 \\
% Max skewness & 11 & 11 \\
% \hline
% \end{tabular}
% \end{table}
\section{Boundary Conditions and Solver Settings}
The initial boundary condition files for both turbulence models were taken from the official OpenFOAM \texttt{motorBike} tutorial, which is included in the standard OpenFOAM installation. This tutorial provides a complete and tested setup for incompressible turbulent flow around a solid body using \texttt{simpleFoam}. After copying these files, the boundary conditions were adapted to match the delta wing simulation, including the correct freestream velocity, angle of attack, domain patch names, and turbulence parameter values.

For the velocity and turbulence variable fields at the outer boundaries of the domain (inlet, outlet, upper wall, lower wall, and side wall), the \texttt{freestream} boundary condition type was used. This type automatically switches between inlet and outlet behaviour depending on the local flow direction at each face. At low angles of attack, the flow enters mainly through the front patch and leaves through the rear. However, at high angles of attack, part of the flow may also enter or leave through the upper, lower, or side boundaries. A fixed-velocity condition on these patches would force the flow in an incorrect direction, while a pure zero-gradient condition could cause instability. The \texttt{freestream} type avoids both problems by applying the prescribed freestream value where the flow enters the domain and a zero-gradient condition where it exits. This makes the setup robust across the entire angle-of-attack range without requiring changes for each case.

For the pressure field, the \texttt{freestreamPressure} type was used on the same outer boundaries. This condition adjusts the pressure treatment depending on whether the flow is entering or leaving the domain at each face. At the outlet patch, a fixed reference pressure value is set, which is required for incompressible solvers to define the absolute pressure level. The symmetry plane at the model centreline uses a \texttt{symmetry} boundary condition to simulate only half of the domain.

On the aircraft surface, a no-slip wall condition is applied for velocity. For pressure, a zero-gradient condition is used. For the turbulence variables, wall-function boundary conditions are applied on the aircraft patch. Wall functions allow the first mesh cell to be placed further from the wall than would be needed if the boundary layer were fully resolved. This significantly reduces the total cell count and computational time while still providing a reasonable estimate of the wall shear stress and turbulent quantities. The choice of wall functions must be consistent with the near-wall $y^{+}$ value discussed in Section~3.3.

Two turbulence models were used: the Spalart--Allmaras (SA) model and the $k$--$\omega$ SST model. These were selected because both were also used in the reference ANSYS Fluent simulations by Giannelis et~al.\ \cite{Giannelis2023SSAMGen5}, which allows a direct comparison. Because each model solves different transport equations, the OpenFOAM case requires different files in the initial conditions directory (the \texttt{0/} folder) and the \texttt{system} folder. The SA model requires a single turbulence field file called \texttt{nuTilda}, while the SST model requires two files \texttt{k} and \texttt{omega} instead. All other files, including velocity (\texttt{U}), pressure (\texttt{p}), and turbulent viscosity (\texttt{nut}), remain the same between the two setups. Switching between the models requires replacing these turbulence variable files and changing the model selection in the \texttt{turbulenceProperties} configuration file.

The simulations were run using the \texttt{simpleFoam} solver with second-order accurate numerical schemes for both convective and diffusive terms. Convergence was monitored by tracking residuals and aerodynamic force coefficients. A simulation was considered converged when the residuals dropped below a set threshold and the force coefficients no longer changed between iterations. The same solver settings, numerical schemes, and convergence criteria were kept identical between the SA and SST simulations to ensure that any differences in the results arise only from the turbulence model.

The complete case configuration, including all boundary condition files, solver settings, and numerical schemes, is provided in Annex~[X].

\section{Post-Processing Methods}


For every iteration the aerodynamic force coefficients were extracted using the \texttt{forceCoeffs} function object available in OpenFOAM. The \texttt{forceCoeffs} utility computes the lift coefficient ($C_L$), drag coefficient ($C_D$), and pitching moment coefficient ($C_M$) by integrating the pressure and viscous (shear) contributions over the selected aircraft surface patches. The output is written during the run, which allows the coefficient histories to be monitored directly during the iteration process \cite{OpenFOAMForceCoeffs,OpenFOAMForces}.

To obtain physically meaningful and comparable coefficients, the required reference values were specified in the \texttt{forceCoeffs} setup: freestream velocity magnitude, reference area, reference length (mean aerodynamic chord), and the moment reference point. These values were selected to match the SSAM-Gen5 reference parameters reported by Giannelis et al.~\cite{Giannelis2023SSAMGen5}, so that the resulting coefficients can be compared directly with both the wind tunnel measurements and the ANSYS Fluent results.

The lift and drag directions were defined explicitly in the \texttt{forceCoeffs} configuration. Since the simulations cover a wide range of angles of attack, the lift and drag directions change for each case. Drag was defined parallel to the freestream direction, while lift was defined perpendicular to the freestream direction. Therefore, for each angle of attack, the corresponding lift and drag direction vectors were recalculated and updated before running the case.

Convergence of the aerodynamic coefficients was assessed by plotting the coefficient histories as a function of iteration number. In well-converged cases, $C_L$, $C_D$, and $C_M$ approach steady values and no longer show a systematic drift. For some high angle-of-attack cases, small oscillations remained even after the residuals decreased, which is typical for highly separated flows. In these cases, the final reported coefficients were obtained by averaging over the last several hundred iterations after the oscillations became statistically stable.

Flow visualisation and qualitative inspection were performed using the open-source software ParaView, which can read OpenFOAM case data directly \cite{Ahrens2005ParaView}. ParaView was used mainly to examine pressure (and pressure coefficient) contours on the delta-wing surface and to visualise flow structures in the domain at different angles of attack. These visual checks help confirm that the solution contains physically reasonable features, such as the development of leading-edge vortices at moderate and high angles of attack.

The near-wall resolution was also verified by checking the $y^{+}$ values on the aircraft surface after each simulation. OpenFOAM provides a dedicated \texttt{yPlus} function object that computes and writes the $y^{+}$ field for turbulence models \cite{OpenFOAMyPlus}. The resulting $y^{+}$ distribution was visualised in ParaView to confirm that the values remain consistent with the chosen near-wall treatment and boundary condition setup across the full surface.

All post-processing steps (coefficient extraction, convergence monitoring, flow visualisation, and $y^{+}$ verification) were applied in the same manner for both turbulence models (SA and $k$--$\omega$ SST) and for all angles of attack. This ensures that the comparisons between turbulence models and against the reference data are consistent.

